\documentclass{article}
\usepackage{graphicx}
\setlength{\parindent}{0pt}

% \title{The Emergence of Axiomatic Attention Patterns During LoRA Fine Tuning of Rerankers}
\title{The Emergence of Relevance Through Axiomatic Attention Patterns During LoRA Fine Tuning}
\author{Matthew Perlman, Atharva Nijasure, James Allan}
\date{April 2026}

\begin{document}

\maketitle

\section*{Abstract}

LoRA and FT are industry standards for domain specific language tasks, and see consistent usage in the training of rerankers.
While its established that FT the attention layers of rerankers is crucial, it is unclear what behavior is learned and where.
Through various ablation and attention mass experiments, we provide a mechanistic interpretation of where and how learned relevance heads emerge, uncovering strong correlatation with axiomatic IR concepts like semantic matching, term rarity, and cross attention.
We find these learned relevance signals are almost exclusive to a small set of mid-network layers, indicating the possibility of more efficient LoRA through sparse FT.

\section*{Introduction}

What is IR and axiomatic ideology, transformers for reranking, FT/LoRA for domain specific models, sparsity of useful heads / lottery ticket hypothesis, so where/what is learned during FT for rerankers.

\vspace{1em}

To determine where and what is learned during FT, we must find where LoRA actually improved performance and what types of features are emerging in these relevance heads.
We therefore perform an initial ablation experiment to identify relevance heads/layers learned during FT, followed by an exploration of how these heads attend to different features before and after FT.

\newpage

\section*{Ablation Experiments}

TODO: describe what we are doing and why.

\subsection*{Methodology}

TODO: explain and justify methodology.

\subsection*{Results}

TODO: what did we uncover, what graphs, quantitative, qualitative results.

\newpage

\section*{Feature Attention Mass Experiments}

TODO: describe what we are doing and why.

\subsection*{Methodology}

TODO: explain and justify methodology.

\subsection*{Results}

TODO: what did we uncover, what graphs, quantitative, qualitative results.

\newpage

\section*{Discussion}

TODO

Combining the results of both experiments, what did we learn. Talk about the sparsity of actually useful LoRA and how such layers are strongly correlated with axiomatic attention patterns.

Why might learning be so sparse? Early layers seem to track with conventional wisdom, but in many other domains FT becomes domain specific in late layers not mid layers. Why is there a DROP in learned relevance for layers 20-25?

Are there features that corresponded more or less? Any unexpected ones? Are these results sufficiently interpretable. Do they meet the threshold for a causal argument?

\section*{Future Work}

TODO

Extensions of the ablation experiment: causal patching, circuit discovery.

Extensions of the attention mass experiment: different features, automatic feature detection.

Generic extensions: Predictive fine tuning (both static and pre-tuning feature presence), possibility of denoising, actually FT manually to demonstrate the result is useful.

\end{document}
